\documentclass[11pt,a4paper]{article}

% --- Encoding ---
\usepackage[utf8]{inputenc}
\usepackage[T1]{fontenc}
\usepackage[spanish,es-noshorthands]{babel}
\usepackage{lmodern}

% --- Layout ---
\usepackage[margin=2.5cm]{geometry}
\usepackage{parskip}
\setlength{\parindent}{0pt}
\setlength{\parskip}{0.5em plus 0.2em minus 0.1em}
\setlength{\emergencystretch}{3em}

% --- Tablas y gráficos ---
\usepackage{booktabs}
\usepackage{array}
\usepackage{enumitem}
\usepackage{xcolor}
\usepackage{graphicx}
\usepackage{amsmath,amssymb}

% --- Soporte para caracteres Unicode en texto ---
\usepackage{newunicodechar}
\newunicodechar{→}{\ensuremath{\rightarrow}}
\newunicodechar{←}{\ensuremath{\leftarrow}}
\newunicodechar{≥}{\ensuremath{\geq}}
\newunicodechar{≤}{\ensuremath{\leq}}
\newunicodechar{×}{\ensuremath{\times}}

% --- Code listings ---
\usepackage{listings}
\definecolor{codegray}{rgb}{0.96,0.96,0.96}
\definecolor{codeborder}{rgb}{0.85,0.85,0.85}
\definecolor{codecomment}{rgb}{0.30,0.55,0.30}
\definecolor{codekeyword}{rgb}{0.10,0.30,0.70}
\definecolor{codestring}{rgb}{0.65,0.15,0.10}

\lstdefinestyle{cstyle}{
    language=C,
    backgroundcolor=\color{codegray},
    basicstyle=\ttfamily\footnotesize,
    breaklines=true,
    frame=single,
    rulecolor=\color{codeborder},
    showstringspaces=false,
    keywordstyle=\color{codekeyword}\bfseries,
    commentstyle=\color{codecomment}\itshape,
    stringstyle=\color{codestring},
    aboveskip=1em,
    belowskip=1em,
}

\lstdefinestyle{shellstyle}{
    language=bash,
    backgroundcolor=\color{codegray},
    basicstyle=\ttfamily\small,
    breaklines=true,
    frame=single,
    rulecolor=\color{codeborder},
    showstringspaces=false,
    keywordstyle=\color{codekeyword}\bfseries,
    commentstyle=\color{codecomment}\itshape,
    stringstyle=\color{codestring},
    aboveskip=1em,
    belowskip=1em,
    columns=fullflexible,
}

\lstdefinestyle{textstyle}{
    backgroundcolor=\color{codegray},
    basicstyle=\ttfamily\footnotesize,
    breaklines=true,
    frame=single,
    rulecolor=\color{codeborder},
    showstringspaces=false,
    columns=fullflexible,
    aboveskip=1em,
    belowskip=1em,
}

\usepackage[hidelinks]{hyperref}

\title{
    \vspace{-2.5em}
    \textbf{Taller 6} \\[0.3em]
    \large Procesos de usuario en EDOS
}
\author{
    Tomás Rodeghiero \\
    \small Sistemas Operativos --- Universidad Nacional de Río Cuarto
}
\date{2026}

\begin{document}

\maketitle

En este taller trabajé sobre el paso \texttt{08-process} de EDOS, que es donde el
kernel pasa a poder cargar y correr programas de usuario. En este informe voy
resolviendo los nueve ejercicios de la consigna: primero el booteo y la
ejecución del \texttt{init}, después el análisis del allocator de páginas, el
layout de memoria del kernel, la creación de procesos, cómo se compila y enlaza
el código de usuario, el camino desde el scheduler hasta la primera instrucción
del proceso, y por último las dos partes con código propio (provocar un page
fault y agregar el syscall \texttt{time}). Todo lo probé en QEMU
(\texttt{qemu-system-riscv32}) y voy pegando las salidas que me sirvieron para
verificar cada punto.

\tableofcontents
\bigskip

% =====================================================================
\section{Booteo de EDOS y ejecución del init}
% =====================================================================

Compilé el paso \texttt{08-process} tal cual viene y lo lancé en QEMU con:

\begin{lstlisting}[style=shellstyle]
qemu-system-riscv32 -machine virt -bios none -nographic \
                    -smp 1 -kernel kernel
\end{lstlisting}

La salida del booteo, ya con paginación habilitada y el primer proceso de
usuario corriendo, es la siguiente (ver \texttt{entrega/salida-original.txt}):

\begin{lstlisting}[style=textstyle]
Initializing kernel block allocator...
free pages=32759, free_list=80009000, kend=80009000, mem_end=88000000
ktext_start: 80000000
kdata_start: 80003000
mapping user code va=00000000 to pa=8002e000
Running scheduler on CPU 0
EDOS init process: pid=1!
Init going to sleep for 4 ticks...
pid 1 going to sleep.
Task pid=1 awake.
init awake! Finishing...
Freeing task init resources...
\end{lstlisting}

Lo que pasa, leyendo el log de arriba para abajo:

\begin{itemize}[leftmargin=1.4em]
    \item \texttt{kernel\_main()} en CPU 0 llama a \texttt{init\_kalloc()}
    (\texttt{kalloc.c}). El log \texttt{free pages=32759} coincide con
    $(128\,\text{MiB} - \texttt{kend}) / 4\,\text{KiB}$ y confirma que la free
    list quedó cargada.
    \item \texttt{init\_vm()} arma el page table del kernel (\texttt{vm.c},
    \texttt{arch.c}). \texttt{map\_kernel\_memory()} imprime
    \texttt{ktext\_start} y \texttt{kdata\_start}.
    \item \texttt{create\_process("init")} en \texttt{task.c} crea la tarea de
    usuario, carga el binario \texttt{init} desde efs y mapea la página del
    código en va=0 → pa=0x8002e000.
    \item El scheduler arranca, selecciona el \texttt{init}, entra a
    \texttt{start\_process()} y vuelve a modo usuario.
    \item \texttt{main()} del init imprime \texttt{"EDOS init process: pid=1!"}
    usando los syscalls \texttt{getpid()} y \texttt{console\_puts()}.
    \item Después llama \texttt{sleep(4)}: el syscall pone al proceso en
    \texttt{WAITING} sobre la condición \texttt{\&ticks}; el log
    \texttt{"pid 1 going to sleep."} sale desde el handler \texttt{sys\_sleep}
    en el kernel.
    \item Cuatro ticks después, \texttt{inc\_ticks()} en \texttt{task.c}
    despierta al proceso (\texttt{"Task pid=1 awake."}) y el init imprime
    \texttt{"init awake!"} y retorna de \texttt{main()}, lo que ejecuta
    \texttt{exit()} en \texttt{edoslib.c}.
    \item El scheduler ve la tarea \texttt{ZOMBIE} y \texttt{free\_resources()}
    la libera.
\end{itemize}

% =====================================================================
\section{Funcionamiento del physical pages allocator}
% =====================================================================

El allocator vive en \texttt{kalloc.c}/\texttt{kalloc.h} y administra el rango de
RAM \texttt{[\_\_kernel\_end, \_\_mem\_end)}, o sea desde donde termina la imagen
del kernel hasta el tope de los 128 MiB de QEMU. La estructura central es una
free list simplemente enlazada cuyos nodos viven dentro de las propias páginas
libres:

\begin{lstlisting}[style=cstyle]
struct link { struct link *next; };
static struct link *free_list = NULL;
static unsigned int free_pages = 0;
\end{lstlisting}

Funciones expuestas (\texttt{kalloc.h}):

\begin{lstlisting}[style=cstyle]
void  init_kalloc(void);
void* alloc_page(void);
void  free_page(void* pa);
\end{lstlisting}

\begin{itemize}[leftmargin=1.4em]
    \item \texttt{init\_kalloc()}: recorre la RAM desde
    \texttt{\_\_mem\_end - PAGE\_SIZE} hacia abajo en pasos de
    \texttt{PAGE\_SIZE} hasta \texttt{\_\_kernel\_end} y llama
    \texttt{free\_page(p)} en cada una. Eso es lo que arma la free list y fija
    \texttt{free\_pages} = cantidad de páginas libres. En el log inicial apareció
    \texttt{free\_pages=32759}, que son los 128 MiB menos lo que ya estaba
    reservado para texto, datos y stacks del kernel.
    \item \texttt{alloc\_page()}: si la lista está vacía retorna 0. Si no, saca la
    cabeza, hace \texttt{memset(n, 0, PAGE\_SIZE)} y la devuelve. El memset
    cumple dos roles: evita filtrar datos del kernel a la página y deja la zona
    en cero, lo que es cómodo para usarla, por ejemplo, como nodo de page table o
    como página de stack.
    \item \texttt{free\_page(pa)}: valida que el puntero esté alineado a 4 KiB
    (\texttt{page\_aligned}) y dentro del rango
    \texttt{(\_\_kernel\_end, \_\_mem\_end)}. Si alguna falla, \texttt{panic}. Si
    pasa, encadena la página en la cabeza de \texttt{free\_list} e incrementa
    \texttt{free\_pages}.
\end{itemize}

Es un allocator extremadamente simple: O(1) en alloc y free, sin fragmentación
externa (todas las páginas son del mismo tamaño), sin metadata aparte del puntero
embebido en la página libre. La única crítica es que no está protegido por un
spinlock; en este paso solo una CPU lo toca, pero en pasos siguientes habría que
serializarlo.

% =====================================================================
\section{Layout del espacio de direcciones del kernel}
% =====================================================================

\texttt{map\_kernel\_memory()} en \texttt{arch.c} configura mapeos 1:1 (va = pa)
para todo lo que el kernel necesita ver con paginación encendida. Las regiones,
leídas en orden del código:

\begin{lstlisting}[style=cstyle]
map_region(pgtbl, UART,        UART,        PAGE_SIZE,    R|W);
map_region(pgtbl, VIRTIO0,     VIRTIO0,     PAGE_SIZE,    R|W);
map_region(pgtbl, PLIC,        PLIC,        PLIC_SIZE,    R|W);
map_region(pgtbl, ktext_start, ktext_start, ktext_size,   R|X);
map_region(pgtbl, kdata_start, kdata_start, kdata_size,   R|W);
\end{lstlisting}

con \texttt{UART = 0x10000000}, \texttt{VIRTIO0 = 0x10001000},
\texttt{PLIC = 0x0c000000}, \texttt{PLIC\_SIZE = 0x400000},
\texttt{ktext\_start = \_\_kernel\_start = 0x80000000},
\texttt{ktext\_size = \_\_text\_end - \_\_kernel\_start} y
\texttt{kdata\_size = \_\_mem\_end - \_\_text\_end}. Diagrama del layout
resultante:

\begin{lstlisting}[style=textstyle]
Direcciones logicas del kernel (1:1 con fisicas)

+---------------------------+ 0x00000000
|                           |
|      no mapeado           |
|                           |
+---------------------------+ 0x0c000000
|   PLIC (R/W)              |   4 MiB (PLIC_SIZE)
+---------------------------+ 0x0c400000
|                           |
|      no mapeado           |
|                           |
+---------------------------+ 0x10000000
|   UART (R/W)              |   1 pagina (4 KiB)
+---------------------------+ 0x10001000
|   VIRTIO0 (R/W)           |   1 pagina
+---------------------------+ 0x10002000
|                           |
|      no mapeado           |
|                           |
+---------------------------+ 0x80000000   __kernel_start
|   kernel .text  (R/X)     |
+---------------------------+ __text_end  (en mi corrida 0x80003000)
|   .rodata, .data, .bss,   |
|   stacks por CPU y resto  |
|   de la RAM (R/W)         |
+---------------------------+ 0x88000000   __mem_end (128 MiB)
\end{lstlisting}

Notas:

\begin{itemize}[leftmargin=1.4em]
    \item Los MMIO se mapean como R/W para que el kernel pueda escribir en UART,
    VIRTIO y PLIC. Nunca llevan \texttt{PAGE\_X}.
    \item La separación \texttt{ktext} (R/X) vs \texttt{kdata} (R/W) impide que el
    kernel ejecute datos por error.
    \item Como los mapeos son 1:1, una ``dirección lógica'' del kernel se
    interpreta igual con paginación encendida o apagada, lo que simplifica
    \texttt{enable\_paging()} (basta cargar \texttt{satp} con el page table;
    ningún símbolo cambia de dirección).
    \item Cada page table del proceso copia este page table del kernel como base y
    le agrega los mapeos U-mode en \texttt{[0, PROC\_MAX\_VA)}, cosa que se ve en
    \texttt{exec()} (\texttt{task.c}).
\end{itemize}

% =====================================================================
\section{Pasos de create\_process() en task.c}
% =====================================================================

\texttt{create\_process(path)} es el camino completo ``de la nada a una tarea de
usuario RUNNABLE''. En orden:

\begin{enumerate}[leftmargin=1.6em]
    \item \texttt{create\_task(path, start\_process)}:
    \begin{itemize}[leftmargin=1.2em]
        \item Busca el primer slot \texttt{UNUSED} en \texttt{tasks[]} (con lock).
        \item Inicializa el contexto en cero
        (\texttt{memset(\&task->ctx, 0, ...)}).
        \item Pone \texttt{task->ctx.ra = start\_process}. Cuando el scheduler
        haga \texttt{context\_switch} hacia este task por primera vez,
        \texttt{ret} saltará a \texttt{start\_process()}.
        \item Reserva la kstack con \texttt{alloc\_page()} y deja
        \texttt{task->ctx.sp = kstack + PAGE\_SIZE} (tope de la pila kernel).
        \item Asigna \texttt{tid}, deja \texttt{pgtbl=kernel\_pgtbl}
        provisoriamente, \texttt{pid=0}, \texttt{state=CREATED}.
    \end{itemize}
    \item Asignación de pid: \texttt{t->pid = ++last\_pid}. Así el init queda con
    pid 1 (es lo que imprime el \texttt{getpid} en el ejercicio 1).
    \item \texttt{exec(t, path, 0)}: es el corazón de la creación del proceso.
    \begin{enumerate}[label=\alph*),leftmargin=1.4em]
        \item Aloca \texttt{new\_pgtbl} con \texttt{alloc\_page()} y copia encima
        el \texttt{kernel\_pgtbl} (\texttt{memcpy}). De esa forma todo proceso
        ``ve'' los mapeos del kernel para cuando entre en modo supervisor por un
        trap.
        \item \texttt{load\_program(new\_pgtbl, path)}:
        \begin{itemize}[leftmargin=1.2em]
            \item \texttt{efs\_file(path)} devuelve el \texttt{struct file} con el
            binario.
            \item Para cada \texttt{PAGE\_SIZE} del binario aloca una página física
            con \texttt{alloc\_page}, copia los bytes y mapea va → pa con flags
            \texttt{PAGE\_R|W|X|U} (visible y ejecutable desde U-mode). La salida
            del ej.\ 1 muestra esto: \texttt{"mapping user code va=00000000 to
            pa=8002e000"}.
        \end{itemize}
        \item Aloca la página del user-stack y la mapea en
        \texttt{PROC\_MAX\_VA - PAGE\_SIZE} con \texttt{PAGE\_R|W|U}. Así el stack
        usuario queda al tope del espacio de proceso.
        \item \texttt{setup\_main\_args(task, ustack + PAGE\_SIZE, args)}: empuja
        las cadenas y arma el array \texttt{char* argv[]}. Deja \texttt{argc} y
        \texttt{argv} en \texttt{a0}/\texttt{a1} del trap frame (a través de
        \texttt{put\_argument}).
        \item Inicializa el trap frame: \texttt{tf->sp = sp} (top del user-stack
        recién armado), \texttt{tf->pc = 0}. Esa \texttt{pc=0} corresponde al
        entry point \texttt{start()} del usuario (ver ej.\ 5).
        \item Si el task ya era un proceso (re-exec), libera la memoria del page
        table anterior con \texttt{unmap\_region(..., true)} y \texttt{free\_page}
        del page table viejo.
        \item \texttt{task->pgtbl = new\_pgtbl}.
        \item \texttt{task->ctx.sp = (size\_t) tf}. Acá está el truco: cuando el
        scheduler haga \texttt{context\_switch}, el \texttt{sp} del task quedará
        apuntando al trap frame en su kstack, y \texttt{u\_trap\_ret} va a leer
        todos los registros de ahí.
    \end{enumerate}
    \item \texttt{t->state = RUNNABLE}. Recién ahora el scheduler puede elegir la
    tarea. La siguiente vez que el scheduler corra en alguna CPU, va a tomar este
    slot y \texttt{context\_switch} hacia él.
\end{enumerate}

Si algo falla, \texttt{exec()} libera lo allocateado y \texttt{create\_process}
retorna false. En el caso del init, \texttt{kernel\_main()} no chequea el
retorno: si fallara, EDOS quedaría sin procesos y solo correría el scheduler
ocioso.

% =====================================================================
\section{user.ld, entry point y mkefs.sh}
% =====================================================================

\subsection*{(a) Linker script user.ld}

\begin{lstlisting}[style=cstyle]
ENTRY(start)

SECTIONS {
    . = 0;
    .text :   { KEEP(*(.text.start)); *(.text .text.*); }
    .rodata : ALIGN(4) { *(.rodata .rodata.*); }
    .data :   ALIGN(4) { *(.data .data.*); }
    .bss :    ALIGN(4) { *(.bss .bss.* .sbss .sbss.*); }
}
\end{lstlisting}

El linker arranca la imagen en virtual address 0 (es lo que espera
\texttt{exec()} cuando hace \texttt{tf->pc = 0}). El \texttt{KEEP} de
\texttt{.text.start} fuerza que la primera función ubicada en \texttt{.text} sea
la marcada con el atributo correspondiente; en este proyecto eso lo logra la
posición física de \texttt{start()} como primer símbolo de \texttt{edoslib.c}
enlazado primero por el Makefile del user (\texttt{edoslib.o usys.o init.c} en
ese orden). El resultado es que la primera instrucción del \texttt{.text}
efectivo es la del \texttt{start()}.

\subsection*{(b) Entry point}

\texttt{ENTRY(start)} le dice al linker que el símbolo \texttt{start} es el punto
de entrada, pero como el kernel ignora la cabecera ELF (el Makefile hace
\texttt{objcopy -O binary init init.bin} y solo eso queda embebido en
\texttt{efsfiles.c}), lo único que importa es que \texttt{start()} caiga en va=0.
La función vive en \texttt{edoslib.c}:

\begin{lstlisting}[style=cstyle]
void start(void) {
    exit(main());
}
\end{lstlisting}

O sea: cuando \texttt{exec()} arme el trap frame con \texttt{tf->pc = 0} y
\texttt{sret} devuelva el control a U-mode, la CPU ejecuta \texttt{start()},
\texttt{start} llama al \texttt{main()} del programa, y cuando \texttt{main}
retorna, \texttt{exit()} (syscall 0) le pasa el código de salida al kernel.

\subsection*{(c) Qué hace mkefs.sh}

El script (lo invoca el Makefile del user con \texttt{"./mkefs.sh init README"})
arma \texttt{efsfiles.c} de la siguiente forma:

\begin{enumerate}[leftmargin=1.6em]
    \item Crea un archivo temporal con un header y un \texttt{\#include "efs.h"}.
    \item Por cada nombre de archivo recibido, si existe
    \texttt{<nombre>.bin} (caso de los programas) usa \texttt{xxd -i
    <nombre>.bin} para volcar el binario como arreglo C; si no usa \texttt{xxd -i
    <nombre>} (caso del archivo de datos README). \texttt{xxd -i} genera dos
    símbolos: el arreglo de bytes \texttt{<nombre>\_bin} (o \texttt{<nombre>}) y la
    longitud \texttt{<nombre>\_bin\_len}.
    \item Con \texttt{sed}, reemplaza \texttt{unsigned int} por
    \texttt{const unsigned int} en las longitudes, para que el linker las ponga en
    \texttt{.rodata}.
    \item Después abre la tabla \texttt{struct file efs\_files\_table[] = \{} y,
    por cada archivo, emite una fila\\
    \texttt{\{"<nombre>", EFS\_FILE\_PROGRAM, <nombre>\_bin, <nombre>\_bin\_len\}}
    o\\
    \texttt{\{"<nombre>", EFS\_FILE\_DATA, <nombre>, <nombre>\_len\}}
    según haya \texttt{.bin} o no.
    \item Cierra la tabla con un centinela \texttt{\{0, 0, 0, 0\}} y mueve el
    \texttt{efsfiles.c} al directorio del kernel.
\end{enumerate}

En tiempo de compilación del kernel, \texttt{efs\_file(name)} (\texttt{efs.c})
recorre esa tabla por \texttt{strcmp} y le devuelve a \texttt{load\_program()} un
puntero al binario y su largo. Así el kernel no necesita lectura de disco: el
filesystem está en RAM, embebido en su \texttt{.rodata}.

% =====================================================================
\section{Del scheduler a la primera instrucción del usuario}
% =====================================================================

Recorrido completo de la primera vez que el init es elegido:

\begin{enumerate}[leftmargin=1.6em]
    \item \texttt{scheduler()} (\texttt{task.c}) corre en CPU 0. Recorre
    \texttt{tasks[]}, toma el lock, encuentra el slot del init en estado
    \texttt{RUNNABLE}, lo pasa a \texttt{RUNNING}, asigna
    \texttt{next\_task->ticks = QUANTUM}, deja
    \texttt{cpus\_state[cpu].task = next\_task} y llama:
    \begin{lstlisting}[style=cstyle]
context_switch(&cpus_state[cpu].ctx, &next_task->ctx);
    \end{lstlisting}
    \item \texttt{context\_switch} (\texttt{arch.s}) guarda \texttt{ra} y
    \texttt{s0-s11} del scheduler en su propia pila (la pila del scheduler de esa
    CPU) y carga desde \texttt{next\_task->ctx}:
    \begin{lstlisting}[style=textstyle]
sp <- next_task->ctx.sp  (== kstack + PAGE_SIZE)
ra <- next_task->ctx.ra  (== start_process)
    \end{lstlisting}
    Finalmente \texttt{ret} salta a \texttt{start\_process}.
    \item \texttt{start\_process()} (\texttt{task.c}):
    \begin{lstlisting}[style=cstyle]
release(&task->lock);
return_to_user_mode();
    \end{lstlisting}
    Libera el lock que tenía el scheduler y va al retorno a usuario.
    \item \texttt{return\_to\_user\_mode()} (\texttt{trap.c}):
    \begin{lstlisting}[style=cstyle]
disable_interrupts();
set_u_trap_handler();         // stvec <- u_trap
set_kstack(kstack+PAGE_SIZE); // sscratch <- kstack top
set_u_previous_mode();        // sstatus.SPP=0, SPIE=1
set_page_table(task->pgtbl);  // satp <- pa2satp(pgtbl)
u_trap_ret(task_trap_frame_address(task));
    \end{lstlisting}
    El \texttt{task\_trap\_frame\_address} es exactamente el lugar de la kstack
    donde \texttt{exec()} dejó armado el trap frame inicial.
    \item \texttt{u\_trap\_ret} (\texttt{trap.s}):
    \begin{lstlisting}[style=textstyle]
mv  sp, a0           # sp apunta al trap frame
lw  a0, 30*4(sp)
csrw sepc, a0        # sepc <- tf->pc (== 0)
lw  ra .. s11        # restaura los 29 registros guardados
lw  sp, 29*4(sp)     # sp <- tf->sp (top del user-stack)
sret
    \end{lstlisting}
    \item \texttt{sret} pasa a U-mode (porque \texttt{sstatus.SPP=0}), habilita
    interrupciones (\texttt{SPIE=1}) y carga \texttt{pc <- sepc = 0}. Como el
    \texttt{pgtbl} del proceso mapea va=0 al código del init, la primera
    instrucción ejecutada en U-mode es la primera del \texttt{.text} del
    \texttt{init.bin}, que es la primera instrucción de \texttt{start()} (gracias
    a \texttt{user.ld} y al orden de linkeo del \texttt{user/Makefile}).
    \item \texttt{start()} ejecuta \texttt{exit(main())}. Como \texttt{exit()}
    está en \texttt{usys.s} (con \texttt{ecall}), la primera vez que el proceso
    necesite interactuar con el kernel cae en \texttt{u\_trap} por SYSCALL y se
    desencadena el camino del ej.\ 8.
\end{enumerate}

% =====================================================================
\section{Provocar un page fault en init.c}
% =====================================================================

Cambié \texttt{init.c} (ver \texttt{entrega/init-ejercicio7.c}) para que el
proceso intente escribir en una dirección no mapeada del espacio del proceso:

\begin{lstlisting}[style=cstyle]
int main(void) {
    printf("EDOS init process: pid=%d!\n", getpid());

    printf("Intentando escribir en una direccion no mapeada "
           "(0xdeadbeef)...\n");
    volatile int *bad = (int *) 0xdeadbeef;
    *bad = 0x42;

    console_puts("Esta linea no deberia ejecutarse.\n");
    return 0;
}
\end{lstlisting}

Por qué dispara la falta:

\begin{itemize}[leftmargin=1.4em]
    \item El proceso solo tiene mapeadas (con \texttt{PAGE\_U}) las páginas de
    \texttt{[0, tamaño del programa)} y la página del stack en
    \texttt{[PROC\_MAX\_VA - PAGE\_SIZE, PROC\_MAX\_VA)}, donde
    \texttt{PROC\_MAX\_VA} es \texttt{0x80000000}.
    \item \texttt{0xdeadbeef} está fuera de \texttt{PROC\_MAX\_VA}. Incluso si
    fuera menor, no está mapeada para U-mode.
    \item El store dispara una excepción \texttt{STORE\_PAGE\_FAULT}
    (\texttt{scause=0x0F}). El handler \texttt{u\_trap} llama a
    \texttt{user\_trap()}, que cae en la rama de ``page fault. Killing task...''
    y deja al task con \texttt{killed=true}. En la salida del scheduler aparece,
    después, \texttt{free\_resources}.
\end{itemize}

Verificación (extracto de \texttt{entrega/salida-ejercicio7.txt}):

\begin{lstlisting}[style=textstyle]
EDOS init process: pid=1!
Intentando escribir en una direccion no mapeada (0xdeadbeef)...
Task init in CPU 0 page fault. Killing task...
Freeing task init resources...
\end{lstlisting}

La línea ``Esta linea no deberia ejecutarse'' no aparece, como esperaba.

% =====================================================================
\section{Camino completo del syscall getpid()}
% =====================================================================

Paso a paso, desde el \texttt{call()} en el código de usuario hasta el
\texttt{ret} en el wrapper:

\begin{enumerate}[leftmargin=1.6em]
    \item Código de usuario (\texttt{init.c}):
    \texttt{"printf(... getpid() ...)"} se resuelve al símbolo externo
    \texttt{getpid} declarado en \texttt{edoslib.h}.
    \item Stub assembler (\texttt{user/usys.s}):
    \begin{lstlisting}[style=cstyle]
getpid:
    li a7, 1     # SYS_GETPID
    ecall
    ret
    \end{lstlisting}
    Carga el número de syscall en \texttt{a7} y ejecuta \texttt{ecall}. La ABI de
    RISC-V deja eventuales argumentos en \texttt{a0..a6}; \texttt{getpid()} no
    tiene.
    \item \texttt{ecall} genera una excepción con \texttt{scause = 0x8}
    (ENVIRONMENT CALL FROM U-MODE) y \texttt{sepc = pc} del \texttt{ecall}.
    \item \texttt{u\_trap} (\texttt{trap.s}): swap \texttt{sp <-> sscratch} (pasa a
    usar la kstack del proceso), reserva \texttt{31*4} bytes para el
    \texttt{struct trap\_frame} y guarda \texttt{ra, gp, t0-t6, a0-a7, s0-s11},
    el \texttt{sp} del usuario (leído de \texttt{sscratch}) y \texttt{sepc} (en
    \texttt{tf->pc}). Después \texttt{call user\_trap()}.
    \item \texttt{user\_trap} (\texttt{trap.c}):
    \begin{itemize}[leftmargin=1.2em]
        \item \texttt{set\_page\_table(kernel\_pgtbl)}: pasa al page table del
        kernel mientras corre el handler.
        \item \texttt{set\_k\_trap\_handler()}: por si dentro del syscall ocurre
        otro trap (en S-mode).
        \item \texttt{case SYSCALL}:
        \begin{lstlisting}[style=cstyle]
enable_interrupts();
syscall(task);
skip_trap_instruction(task_trap_frame_address(task));
        \end{lstlisting}
        \texttt{skip\_trap\_instruction} hace \texttt{tf->pc += 4}, que es lo que
        dejará el \texttt{sepc} tras \texttt{u\_trap\_ret} apuntando a la
        instrucción siguiente al \texttt{ecall}.
    \end{itemize}
    \item \texttt{syscall()} (\texttt{syscall.c}):
    \begin{itemize}[leftmargin=1.2em]
        \item \texttt{n = syscall\_number(tf) = tf->a7 = 1}.
        \item Dispatcher: \texttt{syscalls\_table[1] = sys\_getpid}.
        \item \texttt{sys\_getpid(task)} retorna \texttt{task->pid} (el pid del
        proceso).
        \item \texttt{syscall\_put\_result(tf, result)} deja el valor en
        \texttt{tf->a0}.
    \end{itemize}
    \item \texttt{user\_trap} termina llamando a
    \texttt{return\_to\_user\_mode()} al final (y un \texttt{panic} de
    seguridad por si esa función retornara), así que el flujo no vuelve por
    donde vino sino que rearma el retorno a usuario.
    \item \texttt{return\_to\_user\_mode()} rearma el entorno
    (\texttt{set\_u\_trap\_handler}, \texttt{set\_kstack},
    \texttt{set\_u\_previous\_mode}, \texttt{set\_page\_table} al \texttt{pgtbl}
    del proceso) y llama a \texttt{u\_trap\_ret()} pasándole la dirección del
    trap frame (\texttt{task\_trap\_\allowbreak frame\_\allowbreak address}).
    \item \texttt{u\_trap\_ret} (\texttt{trap.s}) restaura todos los registros
    (incluyendo \texttt{a0} con el resultado del syscall), pone
    \texttt{sepc <- tf->pc} (que ya quedó apuntando a la instrucción siguiente al
    \texttt{ecall}), y \texttt{sret} pasa de nuevo a U-mode.
    \item La CPU ejecuta la instrucción siguiente al \texttt{ecall}, que es el
    \texttt{ret} de \texttt{getpid}. \texttt{a0} ya tiene el pid, así que el
    llamador del usuario recibe el valor de retorno como una función C común.
\end{enumerate}

% =====================================================================
\section{Agregar el syscall time()}
% =====================================================================

Decidí implementar \texttt{time()} devolviendo el contador global de ticks del
kernel (\texttt{volatile unsigned int ticks} en \texttt{task.c}). Es lo natural:
ese \texttt{ticks} lo incrementa \texttt{inc\_ticks()} en cada timer interrupt y
ya hay una función \texttt{get\_ticks()} que lo lee tomando el
\texttt{ticks\_lock}.

Cambios realizados (todos en \texttt{entrega/}):

\textbf{1) syscall.h}
\begin{lstlisting}[style=cstyle]
#define SYS_TIME    6
#define SYSCALLS    7
\end{lstlisting}

\textbf{2) syscall.c} --- agregué \texttt{sys\_time} y lo sumé al final de
\texttt{syscalls\_table[]} (que ahora tiene 7 entradas, igual a
\texttt{SYSCALLS}):
\begin{lstlisting}[style=cstyle]
int sys_time(struct task *task) {
    (void) task;
    return (int) get_ticks();
}
\end{lstlisting}

\textbf{3) user/edoslib.h}
\begin{lstlisting}[style=cstyle]
extern int time(void);
\end{lstlisting}

\textbf{4) user/usys.s}
\begin{lstlisting}[style=cstyle]
.global time
time:
    li a7, 6     # SYS_TIME
    ecall
    ret
\end{lstlisting}

\textbf{5) user/init.c} --- modifiqué \texttt{main()} para imprimir el valor de
\texttt{time()} varias veces, durmiendo 2 ticks entre llamadas:
\begin{lstlisting}[style=cstyle]
for (int i = 0; i < 5; i++) {
    printf("init: time = %d ticks\n", time());
    sleep(2);
}
printf("init: time final = %d ticks. Finishing...\n", time());
\end{lstlisting}

El uso de \texttt{get\_ticks()} (con su spinlock \texttt{ticks\_lock}) en lugar de
leer \texttt{ticks} directamente evita una race con \texttt{inc\_ticks()} en el
caso multi-CPU. Devuelvo \texttt{int} para mantener la firma compatible con el
dispatcher (\texttt{typedef int (*syscall\_f)(struct task *)}).

Verificación (extracto de \texttt{entrega/salida-ejercicio9.txt}):

\begin{lstlisting}[style=textstyle]
EDOS init process: pid=1!
init: time = 0 ticks
pid 1 going to sleep.
Task pid=1 awake.
init: time = 2 ticks
pid 1 going to sleep.
Task pid=1 awake.
init: time = 4 ticks
pid 1 going to sleep.
Task pid=1 awake.
init: time = 6 ticks
pid 1 going to sleep.
Task pid=1 awake.
init: time = 8 ticks
pid 1 going to sleep.
Task pid=1 awake.
init: time final = 10 ticks. Finishing...
Freeing task init resources...
\end{lstlisting}

El contador sube de a 2 entre cada print porque entre llamada y llamada el
proceso duerme 2 ticks; el último \texttt{time()} lee 10 porque hubo 5
iteraciones de \texttt{sleep(2)} además del primer print. Eso valida que el
syscall está bien numerado, el stub pasa por la tabla de syscalls y el resultado
viaja de \texttt{a0} hasta el llamador en user.

\end{document}
